\documentclass[9pt,a4paper,twocolumn,twoside]{rho-class/rho}

% --- 1. Idioma y matemáticas básicas ---
\usepackage[spanish,es-nodecimaldot,es-noindentfirst,es-noshorthands]{babel}


% --- 2. EL TRUCO DE MAGIA (BIEN ESCRITO) ---
% Borramos los comandos problemáticos usando su nombre exacto interno
\expandafter\let\csname not=\endcsname\relax
\expandafter\let\csname not<\endcsname\relax
\expandafter\let\csname not>\endcsname\relax

% Ahora ya podemos cargar la fuente moderna sin que explote
\usepackage{fontspec}
\usepackage{unicode-math}

% --- 3. Bibliografía ---
\addbibresource{rho.bib}

% --- 4. Datos del Informe ---
\title{Estudio de la carga del electrón en el experimento de Millikan}
\author{Edgar Fernández Vigil}
\affil{Física - Facultad de Ciencias - Universidad de Oviedo}
\dates{\today}
\journalname{TEIII - Práctica de Laboratorio}
\journal{Carga del electrón en el experimento de Millikan}

% --- 5. Resumen ---
\begin{abstract}
En este informe se presenta la determinación experimental de la carga elemental $e$ mediante el método de la gota de aceite de Millikan. %
El experimento se fundamenta en el estudio de la dinámica de pequeñas gotas de aceite atomizadas que, tras adquirir carga por fricción, se desplazan en el seno de un campo eléctrico uniforme generado por un condensador de placas. %
Mediante la medición de las velocidades terminales de ascenso y descenso (o suspensión), y considerando el balance entre las fuerzas gravitatoria, de empuje, eléctrica y de fricción viscosa, se calculó la carga total de diversas muestras. %
Para asegurar la precisión de los resultados, se aplicó la corrección de Cunningham a la ley de Stokes, necesaria debido al tamaño microscópico de las gotas en relación con el camino libre medio del aire. %
Finalmente, a través de un análisis estadístico de minimización, se obtuvo el valor de la carga elemental y su incertidumbre, contrastándolo con el valor de referencia de $1.602 \cdot 10^{-19}$ C para verificar la cuantización de la carga eléctrica. %
\end{abstract}
% ----------------------------------------------------------
\begin{document}

    \maketitle
    \thispagestyle{firststyle}

\section{Introducción}

\rhostart{E}ntre los años 1910 y 1913, Robert Andrews Millikan logró determinar el valor de la carga elemental con una exactitud hasta ese momento no alcanzada, demostrando así la cuantización de la carga eléctrica. % [cite: 305]
Por este logro obtuvo el premio Nobel de Física en 1923. % [cite: 306]
Su experimento se basa en la medición de la cantidad de carga de gotitas de aceite que ascienden en el aire dentro del campo eléctrico de un condensador de placas, y descienden cuando no hay campo eléctrico. % [cite: 306]
El valor que determinó en su momento, $e=(1.592\pm0.003)\cdot10^{-19}$ C, se desvía apenas un 0.6\% del valor aceptado en la actualidad. % [cite: 307]

El análisis teórico asume que la gotita tiene forma esférica y está sometida a varias fuerzas: la fuerza de su peso ($F_G$), el empuje ascensional en el aire ($F_A$), la fuerza del campo eléctrico ($F_E$) y la fuerza de fricción de Stokes ($F_{R1,2}$). % [cite: 408, 412, 414, 418]
Del balance de estas fuerzas al ascender con campo y al descender sin él, se pueden extraer expresiones para el radio y la carga de la gotita en función de las velocidades de ascenso ($v_1$) y descenso ($v_2$) medidas. % [cite: 422, 424, 425]
Además, al tratarse de radios muy pequeños que están en el orden de magnitud del camino libre medio de las moléculas del aire, es necesario corregir empíricamente la fuerza de fricción de Stokes para asegurar la validez de los cálculos. % [cite: 467]

\section{Objetivos}

La finalidad de esta práctica es: % [cite: 303]

\begin{itemize}
    \item Comprender experimentalmente la medida de la carga elemental mediante el experimento de Millikan. % [cite: 304]
    \item Realizar la medición de la velocidad de las gotas de aceite empleando dos técnicas distintas: el método de ascenso y el método de la suspensión. % [cite: 427, 448]
    \item Aplicar las correcciones microscópicas a la fricción de Stokes para calcular con precisión el radio corregido ($r$) y la carga corregida ($q$) de cada gotita estudiada. % [cite: 468, 469, 471]
    \item Determinar el número entero de cargas elementales ($n$) presentes en cada gota mediante la minimización directa de la función $F(\hat{e})$. % [cite: 489, 490]
    \item Obtener el valor experimental final de la carga elemental ($e$) y su respectiva incertidumbre realizando una media ponderada de las medidas, para finalmente compararlo con el valor de referencia de $1.602\cdot 10^{-19}$ C. % [cite: 304, 496, 497, 500]
\end{itemize}

\section{Procedimiento Experimental}

El procedimiento se basa en el seguimiento óptico de gotas de aceite mediante un microscopio de medición y una cámara acoplada a un ordenador para el registro de tiempos. % [cite: 309, 310]
El desarrollo de la toma de datos se divide en las siguientes fases:

\subsection{Preparación y selección de la muestra}
Se introdujeron gotas de aceite en la cámara de experimentación mediante una compresión corta y fuerte de la bola del soplador conectada al pulverizador. % [cite: 381, 382]
Utilizando la iluminación LED y el microscopio, se seleccionó una gota que descendiera lentamente (aproximadamente entre 0.025 y 0.1 mm/s) y que apareciera como un punto brillante y definido en el foco. % [cite: 384, 393]

\subsection{Método de ascenso y descenso (Dinámico)}
Una vez seleccionada una gota apropiada, se procedió a cronometrar su paso entre las marcas de la escala del ocular (típicamente una distancia de 1 a 2 mm): % [cite: 385, 443]
\begin{enumerate}
    \item \textbf{Descenso:} Con el campo eléctrico desactivado ($U = 0$), se midió el tiempo $t_2$ que tarda la gota en caer por efecto de la gravedad. % [cite: 432, 442]
    \item \textbf{Ascenso:} Se activó la tensión $U$ con la polaridad adecuada para que la gota ascendiera, midiendo el tiempo $t_1$ de tránsito entre las mismas marcas. % [cite: 430, 436]
    \item Este proceso se repitió al menos tres veces por cada gota para minimizar la incertidumbre estadística. % [cite: 385]
\end{enumerate}

\subsection{Método de la suspensión (Estático)}
Como alternativa, se empleó el método de suspensión para determinadas gotas: % [cite: 448]
\begin{enumerate}
    \item Se ajustó la tensión $U$ de las placas hasta que la fuerza eléctrica equilibró exactamente el peso efectivo de la gota, manteniéndola inmóvil en una posición deseada del ocular. % [cite: 453]
    \item Se anotó el valor de la tensión de suspensión $U$ indicada en la unidad de mando. % [cite: 454]
    \item Posteriormente, se desconectó la tensión y se midió el tiempo de descenso $t_2$ en caída libre para determinar el radio de la gota. % [cite: 455, 461]
\end{enumerate}

\subsection{Procesado de datos y correcciones}
Para cada gota medida, se calcularon las velocidades $v_1$ y $v_2$ teniendo en cuenta el aumento del objetivo (2x). % [cite: 443, 475]
A partir de estos valores, se determinó el radio inicial $r_0$ y la carga aparente $q_0$. % [cite: 446, 463]
Debido a que el radio de las gotas es comparable al camino libre medio de las moléculas de aire, se aplicó la corrección de Stokes utilizando la presión atmosférica local para obtener la carga real corregida $q$. % [cite: 467, 471]
Finalmente, se utilizó un algoritmo de búsqueda del entero $n$ que minimiza la dispersión de las cargas para extraer el valor de la carga elemental $e$. % [cite: 490, 495]

\section{Resultados y Discusión}
holaaa

\end{document}
